\documentclass[a4paper]{article}

\usepackage{graphicx}
\usepackage{amsmath,amssymb}
\usepackage{minted}
\usepackage{multirow}
\usepackage{float}
\usepackage{todonotes}
\usepackage{forest}
\usepackage{xstring}
\usepackage{xcolor}
\usepackage[most]{tcolorbox}
\usepackage{xparse}
\usepackage{natbib}

\usetikzlibrary{graphs,graphdrawing}
\usegdlibrary{layered}

% for external graphviz
\input{/home/donkarlo/Dropbox/repo/nd_ai_project/src/nd_ai/knoweldge_representation/language/formal/computer/programming/tex/latex/code_clips/external_graphviz.tex}

% for tags
\input{/home/donkarlo/Dropbox/repo/nd_ai_project/src/nd_ai/knoweldge_representation/language/formal/computer/programming/tex/latex/parts/control_sequence/kind/semtag/semtag.tex}

% for links
\input{/home/donkarlo/Dropbox/repo/nd_ai_project/src/nd_ai/knoweldge_representation/language/formal/computer/programming/tex/latex/code_clips/seehere.tex}

\title{}
\date{}

\begin{document}
\maketitle
    Neural computation is the information processing performed by networks of neurons. Neural computation is affiliated with the philosophical tradition known as Computational theory of mind, also referred to as computationalism, which advances the thesis that neural computation explains cognition.
    \url{https://en.wikipedia.org/wiki/Neural_computation}

    \bibliographystyle{plainnat}
    \bibliography{/home/donkarlo/Dropbox/repo/nd_shared_research/bibliography/bibliography.bib}
\end{document}